\section{A governed upgrade constitution for recursive research systems}
\label{sec:upgrade}

The central engineering risk in self-improvement is semantic collapse between four distinct events:
\rakldisp{
\text{code exists}
\neq
\text{code is deployed}
\neq
\text{improvement is supported}
\neq
\text{the method is the governed incumbent}.
}
A software branch can be implemented and even merged for operational reasons while its scientific method claim remains unvalidated. The live v3 deployment supplies an instructive example: the v3 branch was merged into \texttt{main} under an explicit operator request before its pending CI had completed, and later commits repaired package/publication CI. That sequence is preserved as process history. It is not represented as evidence that v3 was superior merely because it was deployed.

\subsection{Method identity is not the package version}

The inspected Python package still reports software version \texttt{0.1.0}. We therefore separate software/API compatibility from research-method identity. A complete experimental citation should bind at least
\rakldisp{
V=(v_{\mathrm{pkg}},v_{\mathrm{method}},e_{\mathrm{constitution}},\Sigma_{\mathrm{schemas}},h_{\mathrm{git}},h_{\mathrm{validators}}).
}
We use \texttt{3.0.0} as the method identity for the deployed v3 architecture at the Paper 5 freeze, while the exact Git subject remains the decisive reproducibility identity. Patch-level method versions such as 3.0.x are reserved for implementation corrections that do not intentionally change research policy. Minor method versions such as 3.1 and 3.2 denote workflow/method challengers. A change to protected authority/governance semantics is constitutional and should advance a separately visible constitution epoch or major method generation rather than being hidden in an ordinary workflow patch.

The Paper 5 branch introduces a proposal-only \texttt{RAKL\_VERSION.json} manifest to make these coordinates machine visible. Its presence does not promote itself to canonical governance.

\subsection{Change classes}

We retain the existing three-way classification in \texttt{meta.py}.

\paragraph{Class A: implementation.}
Examples include parser defects, CI portability, serialization, hashing, deterministic replay and publication build repairs. A Class-A change requires a reproducing regression, relevant invariant tests, artifact identity and preservation of history. It can improve engineering reliability without supporting a method-capability claim.

\paragraph{Class B: workflow or method.}
Examples include retrieval policy, fibre compilation, routing, saturation, gluing, lesson induction, memory consolidation and research-portfolio allocation. A Class-B challenger must state a positive meta-QoI prediction before evaluated outcomes, run matched parent/challenger development tests, include fresh transfer and hostile near-misses, preserve blocking invariants, use comparable resources, and pass protected fresh assurance before a strong evolution claim.

\paragraph{Class C: constitution.}
Examples include changing what evidence may mint authority, weakening review independence, changing the meaning of missing evidence, or changing who can promote an incumbent. Constitutional changes are never auto-promoted. They require explicit human-visible amendment review, migration/rollback analysis and adversarial authority tests.

\subsection{Roles and anti-self-certification}

A strong evolution claim separates at least seven roles, even if one organization operates them:

\begin{enumerate}[leftmargin=*]
\item \textbf{observer} records content-bound research episodes and process telemetry;
\item \textbf{upgrade proposer} maps repeated evidence to a falsifiable framework hypothesis;
\item \textbf{challenger engineer} implements the frozen proposal;
\item \textbf{development evaluator} runs visible debugging/replay tests;
\item \textbf{fresh assurance evaluator} uses a packet frozen before mutation and hidden from the optimizer;
\item \textbf{governance authority} decides whether an assured variant becomes incumbent;
\item \textbf{post-promotion attestor} observes the actual active repository state and exact-main validation.
\end{enumerate}

Same-session role labels do not create independence. Process independence and evidence-lineage independence must be separately established when independent credit is claimed.

\begin{proposition}[No self-promotion from proposal output]
Suppose a challenger cannot modify the protected evaluator, assurance packet, promotion thresholds or governance credentials that judge it, and promotion requires a content-bound certificate issued outside the challenger path. Then the challenger's generated narrative or code output alone is insufficient to establish its own promotion.
\end{proposition}

\begin{proof}
The promotion predicate contains at least one required input whose value is controlled outside the challenger and is bound to the exact challenger identity. Changing the candidate output cannot set that external input to an accepting value. Therefore self-generated output is insufficient for promotion, even if it can affect development metrics. The claim is conditional on the protection and identity assumptions; compromise of those boundaries remains a separate threat model.
\end{proof}

\subsection{Upgrade proposal packet}

Before implementation outcomes are observed, a Class-B/C proposal freezes a machine-readable packet containing the parent method/Git identity, source episode and diagnosis IDs, change class, alternative diagnoses, proposed mechanism, affected method contracts, predicted primary and secondary meta-QoIs, material-effect thresholds, possible regressions, negative controls, hostile near-misses, development benchmark, fresh assurance reserve, model/tool/resource contract and rollback plan. If the result predates the packet, it is retrospective calibration only.

This packet is intentionally stricter than an issue description. It converts the statement ``we think this might help'' into a falsifiable prediction about a content-identified framework mutation.

\subsection{Challenger lifecycle and variant DAG}

We model evolution as a DAG rather than a destructive sequence:
\rakldisp{
\begin{gathered}
\texttt{OBSERVED\_PATHOLOGY}
\rightarrow\texttt{UPGRADE\_HYPOTHESIS}
\rightarrow\texttt{PREREGISTERED}
\rightarrow\texttt{CHALLENGER\_IMPLEMENTED}\\
\rightarrow\texttt{DEVELOPMENT\_VALIDATED}
\rightarrow\texttt{FRESH\_ASSURANCE\_PENDING}
\rightarrow\{\texttt{ASSURED},\texttt{REJECTED},\texttt{META\_OVERFIT},\CANNOTCHECK\}\\
\rightarrow\texttt{GOVERNANCE\_APPROVED}
\rightarrow\texttt{PROMOTED}
\rightarrow\texttt{ACTIVE\_PROMOTION\_ATTESTED}
\rightarrow\texttt{MONITORED}.
\end{gathered}
}
Rejected and superseded variants remain negative history. A previous incumbent remains a rollback target when compatible.

The existing v3 \texttt{EvolutionArchive} already separates challenger, assured, rejected and incumbent states and does not auto-promote a challenger merely because an evolution trial passes. However, its current \texttt{promote\_incumbent} interface accepts a caller-provided \texttt{governance\_approved} Boolean. We regard that as an implementation gap, not as adequate promotion evidence. A future 3.1 challenger should replace declaration-bound governance with a content-bound attestation naming exact candidate, assurance verdict, constitution epoch and approving process.

\subsection{Development evaluation versus fresh assurance}

Two evidence phases answer different questions. Visible development evaluation can guide debugging and optimization; once exposed, it is no longer blind assurance. Strong method-evolution evidence therefore consumes a separate fresh packet frozen before the challenger and hidden from the proposer. Repeated peeking consumes the assurance exposure budget. The existing \texttt{SelfEvolutionAssessor} and \texttt{evaluate\_bootstrap\_trial} already distinguish development gain, fresh transfer, evidence-lineage independence, evaluator separation, resource matching and meta-overfit. These components support a scoped verdict rather than a global ``Orion evolved'' state.

\subsection{Evaluator migration}

A reviewer may reasonably object that fixed evaluators eventually become obsolete. We allow evaluator evolution, but not inside the candidate it judges. A legitimate evaluator change is a separately governed parent mutation. Within an evaluation epoch, the ruler is frozen. At epoch boundaries, a new evaluator can be proposed, benchmarked against known-answer and adversarial cases, versioned and then used prospectively. This differs from allowing the optimizer to rewrite its reward after seeing an inconvenient score.

\subsection{Operator override and deployment before assurance}

Human or repository governance can intentionally override the scientific promotion process. The protocol therefore models an explicit \texttt{OPERATOR\_OVERRIDE} with requested deviation, exact subject, blockers known at the time and containment/rollback plan. An override may move code if repository permissions allow, but creates no evolution-evidence credit. The deployed variant remains method-provisional until the required evidence is produced. This rule is not hypothetical; it is motivated by the actual v3 merge-before-CI incident and prevents the paper from retrospectively rewriting operational history into a clean preregistered story.

\subsection{Post-promotion attestation and rollback}

Pre-promotion authorization is not evidence that deployment occurred correctly. The existing \texttt{promotion\_attestation.py} distinguishes a candidate that passed checks from an active state whose \texttt{main} ref descends from the candidate, required content matches and exact-active-main post-promotion validation succeeds. A promoted variant also carries an exact rollback parent, migration notes, rollback test and monitoring conditions. New evidence can narrow, supersede or roll back a variant without erasing the evidence that originally supported it.

\subsection{Version increments as empirical claims}

Under this constitution, a label such as \texttt{3.1} has scientific content. It means that a distinct workflow/method challenger has been identified and evaluated under the registered process; it does not mean that a certain number of commits were merged. Consequently Paper 5 can plot \RAKL{} 3.0, 3.1 and later versions as experimental interventions rather than arbitrary software milestones.